\scanpage{56}
\mspage{II. 4}
% lines 1-4
alors de \deleted{considérer} \inserted{prendre} l'image inverse $V$ de $U$ dans $\pi$, la restriction de la \unsure{$u$} à $V$ provient d'une représentation de l'image de $V$ dans $\pi_0^{\mathrm{ab}}(\ell)$, donc d'une repr. de l'image de $V$ dans $\pi_0$, cqfd.

\diagram{\begin{tikzcd}[ampersand replacement=\&,column sep=small,row sep=small]
\& V\arrow[r,hook]\arrow[d]\arrow[dl] \& \pi\arrow[d]\arrow[dll]\\
\operatorname{Im}V\subset\pi_0\arrow[dr]\arrow[drr] \& U\arrow[r,hook]\arrow[d,"\unsure{\wr|\simeq}" description] \& \pi^{\mathrm{ab}}(\ell)\arrow[d]\\
\& U\arrow[r,hook] \& \pi_0^{\mathrm{ab}}(\ell)
\end{tikzcd}}

% lines 5-6
On va donner une application de ces résultats à la théorie des V.A.

% lines 7-12, rewritten statement
\uline{Théorème} 1.5. Soient $k$, \deleted{$X,\unsure{\xi|\zeta}$ comme dans 1.1, \unsure{et}} \inserted{un corps sép. clos, $X$ schéma loc. de type fini sur $k$, normal et \deleted{\unsure{géométriquement}} \inserted{géométriquement intègre}, $\xi$ est pt générique de $X$,} \deleted{\unsure{sous} \inserted{en supposant le corps de type fini mais sép. clos}} \deleted{\unsure{que}} $A$ un schéma abélien sur $X$, et \deleted{\unsure{soit}} $\ell$ un nombre premier distinct de la caractéristique de $k$. Supposons que la représentation de $\pi_1(\overline X,\xi)$ dans $T_\ell(A_\xi)$ soit \deleted{\uline{\unsure{semi-abélienne}}} \inserted{essentiellement résoluble.}

\diagram{\begin{tikzcd}[ampersand replacement=\&]
A\arrow[d,dash]\\X\arrow[d,dash]\\S
\end{tikzcd}}

% lines 13-21
Alors cette représentation est essentiellement \deleted{unipotente, donc essentiellement triviale si elle est semi-simple $(*)$. Dans ce \unsure{cas},} \inserted{triviale, et} il existe un revêtement étale fini surjectif $X'\longrightarrow X$, \deleted{et une extension finie \unsure{de type} finie séparable $k'$ de $k$, \unsure{dont} \unsure{un} $k$-morphisme $X'\longrightarrow\operatorname{Spec}k'$ \unsure{fais}} un schéma abélien $B$ sur $k$, \deleted{tel que $A_{X'}$} et une \uline{isogénie} \inserted{\uline{radicielle}} $B_{X'}\mathrel{\text{---}}A_{X'}$. \deleted{qui soit radicielle} En particulier, si $k$ est de car. nulle, \deleted{\unsure{donc}} $A_{X'}$ est isom. à $B_{X'}$ (i.e. est un schéma abélien constant).

% lines 22-25
Prouvons d'abord l'assertion que la représentation de $\pi_1(X,\xi)$ dans $T_\ell(A_\xi)$ est \deleted{unipotente \unsure{en car.}} \inserted{ess. unipotente}. \deleted{Dans} \deleted{\unsure{ceci}} Soit $U$ un voisinage ouvert affine de $\xi$,

% footnote, cancelled
\deleted{$(*)$ \unsure{Il y a tout lieu de conjecturer} que sans hypothèse supplémentaire l'action de $\pi_1(\overline X,\xi)$ sur $T_\ell(A_\xi)$ est \uline{toujours} \unsure{$\frac12$ simple|semi-simple}.}
% archive sheet number: 55
